\FloatBarrier
\section{Magnets duplicate shielding removal}
\label{sec:magnets-dedup}

Upcoming in GEMC2 6.2 and CLAS12 GEMC Systems 0.4, the October 21 magnets cleanup removes nine
\texttt{DSShield*} meshes and apex from the magnets CAD collection because this shielding is owned by the
beamline. All three configurations are updated in the GEMC2 and GEMC3 sources, runtime collections, and SQLite
banks. The retained 51 magnet meshes, their materials and passive placements, the native solenoid, and all
beamline geometry are unchanged. Load the beamline alongside magnets to include the shielding.

This is removal of redundant geometry, without retriangulation or surface repair. The magnets collection falls
from 61 to 51 source meshes, from 148,940 to 132,230 stored triangles (16,710 removed; 11.22\%), and from
7,452,124 to 6,615,784 bytes (836,340 removed; 11.22\%). The nine downstream meshes account for 14,272
triangles and apex for 2,438. Each retained mesh matches its pre-removal bytes in all four source/runtime
collections.

The benchmark retains the native beamline and solenoid, the complete run-11 beamline CAD bank, and all
remaining magnets CAD. The original fixture includes the ten obsolete placements; the optimized fixture
removes them. Both use the same fixed GEMC3 dev executable, released October 19, 2026, Geant4 11.4.3, macOS
arm64, FTFP\_BERT, one thread, and default production cuts. There is no magnetic field, optical physics, GUI,
event output, target, luminosity normalization, or pileup. The 5 GeV/c electron, proton, positive-pion, and
neutron primaries start at the origin, uniformly distributed in solid angle from 0.25 to 3.5 degrees over full
azimuth. Five alternating-order repetitions use 200 events per particle and seeds 12345--12349, with identical
LUND inputs within each seed pair. All 52 processes complete; every tracking process reports 200 events.

Table~\ref{tab:magnets-dedup-inventory} lists the removed source meshes.
\begin{table}[htbp]
\centering
\caption{Redundant magnets meshes removed from all three configurations.}
\label{tab:magnets-dedup-inventory}
\begin{tabular}{lr}
\toprule
Removed mesh & Triangles \\
\midrule
DSShieldBackAss & 988 \\
DSShieldBackCover & 1,012 \\
DSShieldBackLead & 1,244 \\
DSShieldFlange & 5,172 \\
DSShieldFlangeAttachment & 1,888 \\
DSShieldFrontAss & 1,056 \\
DSShieldFrontCover & 1,252 \\
DSShieldFrontLead & 988 \\
DSShieldInnerAss & 672 \\
apex & 2,438 \\
\bottomrule
\end{tabular}
\end{table}

Median zero-event process wall time changes from 1.98732 to 1.89090 s (1.051 times faster), CPU time from
1.920 to 1.840 s, and peak RSS from 321.53 to 306.98 MiB (4.52\% lower). Tracking times below are median
Geant4 beamOn wall times excluding initialization; speedup is before/after. These short, isolated workloads do
not establish a full-CLAS12 gain. Removing intersecting material changes transport histories, so timing
changes cannot be attributed solely to fewer facets or taken as proof of response equivalence.

Table~\ref{tab:magnets-dedup-tracking} reports tracking separately from loading.
\begin{table}[htbp]
\centering
\caption{Magnets and beamline tracking; five repetitions of 200 events per particle.}
\label{tab:magnets-dedup-tracking}
\begin{tabular}{lrrrr}
\toprule
Particle & Before (s) & After (s) & Speedup & Paired ratio range \\
\midrule
Electron & 35.49510 & 34.50800 & 1.029 & 1.002--1.029 \\
Proton & 10.53440 & 10.34530 & 1.018 & 1.007--1.098 \\
Positive pion & 14.12550 & 14.32390 & 0.986 & 0.981--1.049 \\
Neutron & 10.53630 & 10.27570 & 1.025 & 0.957--1.078 \\
\bottomrule
\end{tabular}
\end{table}

The matched zero-event construction diagnostics use check\_overlaps=2 and seed 12345. Built volume counts,
including the world, fall from 95 to 85; GeomVol1002 exception blocks fall from 58 to 26. Each fixture still
reports 26 GeomSolids1001 and 11 GeomSolids1002 blocks per import. Warning counts are sampled diagnostics, not
counts of all intersecting pairs, and the retained invalid/rejected-facet warnings and assembly contacts
remain unresolved. All three updated runs pass standalone GEMC3 magnets construction (53 volumes including the
world) and combined magnets/beamline GEMC2 TEXT/CAD and SQLite/CAD initialization. The GEMC3 semantic
comparator matches one native solenoid row and 51 CAD rows per configuration; independent checks match the
effective GEMC2 SQLite CAD rows to the GXML placements and the native TEXT rows to SQLite. All unrelated GEMC2
CAD rows and beamline meshes are unchanged. A supplemental combined GEMC3 card constructs 85 volumes at runs
11 and 3029, but only 53 at run 4763: the current GEMC3 lookup omits the beamline banks without that exact run
key, where GEMC2 selects the earlier run-3029 banks. This existing run-selection limitation is outside the
mesh-removal change and remains unresolved. The timed run-11 comparison loads both systems and is unaffected.
Across the forty timed tracking processes, navigation exception blocks total GeomNav1002: 4 before and 2
after.

The dedicated data/magnets\_dedup directory records all measurements, mesh hashes, removed-facet inventory,
database checks, and validation. magnets\_dedup\_archive.zip retains the complete before/after fixtures,
original collected inputs, all raw logs and process reports, matched primary files, source snapshots, and
reproduction scripts. The executable is identified by hash and is not redistributed. Earlier study archives
remain unchanged.
